\section{Introduction}
\IEEEPARstart{R}{obotic} baggage handling is becoming an important component of intelligent airport logistics, where sorting, loading, and transfer operations must be performed reliably within strict time windows. Vision feedback is essential in such systems because it connects object perception with grasp planning and execution. Classical visual servoing establishes closed-loop control by mapping visual feature errors to robot motion commands \cite{hutchinson2002tutorial}. However, in practical baggage handling environments, the dominant bottleneck is often not the control law itself, but the reliability and geometric fidelity of the visual observations used by the controller. Baggage items are visually heterogeneous, frequently occluded, and often equipped with deformable tags, straps, wheels, and handles. These factors make the perception-to-action pipeline substantially more difficult than standard pick-and-place manipulation with regular objects.

Conventional robotic grasping pipelines usually detect the target in RGB images and then fuse depth information for 6D pose or 3D bounding-box estimation. Bounding-box-based detection is efficient but lacks the boundary-level accuracy required for precise manipulation, especially when objects are stacked, partially occluded, or densely arranged. Instance segmentation improves spatial localization by separating objects at the pixel level, but segmentation errors at the boundary may still be amplified after projection into the point cloud. In baggage scenarios, flexible tags and straps are particularly problematic: they may be segmented as part of the rigid body in the image, but they do not belong to the stable geometric structure needed for grasping and placement.

Airport baggage handling further imposes several coupled challenges. First, visual interference arises from highly salient and deformable components such as identification tags, ropes, ribbons, and straps. General purpose segmentation models \cite{he2017mask} may regard these components either as independent objects or as part of the suitcase body, which shifts the estimated centroid and degrades grasp accuracy. Second, 3D projection distortion occurs when residual 2D mask leakage produces outliers in the projected point cloud. Such outliers can bias PCA-based pose estimation and OBB construction \cite{rusu20113d, gottschalk1996obbtree}. Third, temporal instability is caused by conveyor motion, sliding, rolling, collision-induced disturbances, motion blur, and depth-sensor noise. These factors introduce frame-to-frame jitter that directly affects predictive grasping. Finally, high-tempo industrial constraints require the robot to grasp moving targets within a short time window. In our system, the conveyor speed reaches approximately \revadd{0.5~$\mathrm{m/s}$}. Downstream loading also requires high space utilization, so even small pose errors or volume-estimation errors may reduce packing density and increase the number of required loading cycles.

Suction grasping is attractive for dense baggage scenes because it requires limited lateral insertion space, unlike a mechanical gripper that often needs side clearance. Nevertheless, suction grasping is sensitive to local surface localization and pose accuracy. A successful suction grasp therefore depends on reliable segmentation, stable 3D modeling, and accurate motion prediction \cite{bohg2013data}. In this work, a dedicated hybrid end-effector is designed that integrates both suction and a gripper mechanism. \revadd{The gripper is the primary mode for normal dynamic grasping. Suction is used for coarse relocation during active state regulation when confidence in the target grasp is low, adjusting the object pose before state re-estimation and grasp-trajectory replanning.} This hardware configuration motivates a closed-loop software architecture in which \revadd{instance segmentation, spatio-temporal geometric modeling, motion planning, and active state regulation are coordinated.}

To address the above issues, this paper proposes a vision-driven closed-loop perception, geometric modeling and control framework for high-precision robotic baggage handling. At the perception level, a Multi-Scale Feature Aggregation and Attention-Guided Mask Optimization Network (MSAM-Net) is developed \revadd{using a YOLOv10 backbone} to obtain accurate instance masks under tag and strap interference. At the geometric modeling level, a pixel-to-3D Cross Consistency Conditional Filtering (PX3C) method removes residual flexible outliers after 2D-to-3D projection \revadd{within a PCL pipeline with PCA-based OBB estimation}, and a Consistency-Aware OBB Stabilization (CAOS) strategy enforces temporal smoothness across frames. At the action level, a discrete Kalman-filter-based recursive Bayesian estimator predicts target states, a two-stage relative-velocity planning strategy executes rapid interception followed by precise synchronization, and an active state-regulation mechanism resets highly uncertain baggage configurations when passive observation is unreliable.

The contributions of this paper are summarized as follows.
\begin{itemize}
  \item A real-time semantic perception network, MSAM-Net, is proposed for baggage instance segmentation under severe visual interference. The network combines four-level multi-branch feature aggregation with attention-guided multi-scale mask refinement.
  \item A spatio-temporal geometric modeling framework is introduced. PX3C removes residual 3D interference by enforcing pixel-to-3D geometric consistency, and CAOS stabilizes OBB pose and scale estimates using quality-aware temporal smoothing.
  \item A closed-loop dynamic grasping strategy is developed by integrating recursive Bayesian state estimation, two-stage relative-velocity planning, and active state regulation. The complete system is validated on a real robotic baggage-handling platform.
\end{itemize}